\documentclass[11pt]{article}
\usepackage[margin=1in]{geometry}
\usepackage{setspace}
\usepackage{enumitem}
\usepackage{amsmath}
\usepackage{amssymb}
\usepackage{fancyhdr}
\IfFileExists{lastpage.sty}{\usepackage{lastpage}}{}
\usepackage{hyperref}

\setlength{\parindent}{0pt}
\setlength{\parskip}{0.6em}
\setlist[itemize]{leftmargin=*}
\setlist[enumerate]{leftmargin=*}

\newcommand{\answerlines}[1]{\vspace{#1\baselineskip}}

\pagestyle{fancy}
\fancyhf{}
\IfFileExists{lastpage.sty}{
  \fancyfoot[C]{Financial Data Science II\ \ \ Updated March 18, 2026\ \ \ Page \thepage\ of \pageref{LastPage}}
}{
  \fancyfoot[C]{Financial Data Science II\ \ \ Updated March 18, 2026\ \ \ Page \thepage}
}

\begin{document}

\begin{center}
{\Large \textbf{Part 25: Factor Exposures and Risk Adjustment}}
\end{center}

A selected feature or signal may still be difficult to interpret if it is
largely a disguised exposure to common risk factors. For that reason, it is
often useful to ask how much of a signal's behavior is explained by broad
factor movements and how much remains after risk adjustment.

In this Part we use a simple rolling factor model to illustrate two ideas:

\begin{enumerate}
  \item estimating factor exposures
  \item constructing a residual, or risk-adjusted, version of a signal
\end{enumerate}

\textbf{Exercise.} Why might a long-short strategy with strong raw returns
still be disappointing if most of its performance comes from broad market
exposure?

\answerlines{6}

\newpage

\section*{Basic Terms}

A \textbf{factor exposure} is the sensitivity of an asset or strategy to a
common source of variation. A \textbf{beta} is a regression coefficient
measuring this sensitivity. \textbf{Alpha} is the component of return not
explained by the factor specification being used. A \textbf{residual return} is
the realized return minus the fitted value from a factor model. A
\textbf{risk-adjusted signal} is a signal formed after removing some common
factor component.

None of these quantities is model-free. They depend on the chosen factor set
and estimation window.

\textbf{Exercise.} Explain why the same stock can have different estimated
factor exposures under different factor specifications.

\answerlines{6}

\newpage

\section*{Why Risk Adjustment Matters}

Risk adjustment is useful for at least three reasons.

\begin{enumerate}
  \item It helps distinguish common factor exposure from asset-specific information.
  \item It provides a cleaner basis for comparing raw and residual signals.
  \item It helps explain why a strategy behaves as it does during broad market moves.
\end{enumerate}

In practice one rarely removes all systematic variation. The point is to
understand what the signal is loading on and whether those exposures are
desired.

\section*{An Illustrative Factor Set}

For illustration we use the following proxy factors:

\begin{enumerate}
  \item \texttt{SPY} for the broad market
  \item \texttt{IWM - SPY} as a size spread proxy
  \item \texttt{IWD - IWF} as a value-versus-growth proxy
\end{enumerate}

These are only rough proxies, but they are sufficient for an instructional
example.

\newpage

\section*{Rolling Exposures and Residual Returns}

A rolling factor model produces two useful objects:

\begin{enumerate}
  \item estimated exposures through time
  \item residual returns after removing the fitted factor component
\end{enumerate}

These may then be used to compare a raw momentum signal to a momentum signal
built from residual returns.

\section*{Interpreting the Comparison}

If the residual portfolio has smaller factor loadings than the raw portfolio,
then the adjustment has succeeded in reducing common-factor exposure under the
chosen model. This does not imply that the residual portfolio is superior. It
may have lower return, higher turnover, or less stable behavior.

The main purpose of the exercise is diagnostic. It helps answer the question:
is the signal primarily a repackaged factor bet, or does it retain some
distinct component after adjustment?

\newpage

\section*{Exercises}

\textbf{Exercise 1.} Change the rolling estimation window from 24 months to 36
months. How do the exposure estimates change?

\answerlines{6}

\textbf{Exercise 2.} Replace the factor proxies with a different set of ETFs.
Which results are most sensitive to this change?

\answerlines{6}

\textbf{Exercise 3.} Why is it useful to compare the factor exposures of the
portfolio, rather than only the factor exposures of the individual stocks?

\answerlines{6}

\textbf{Exercise 4.} Give an example of a signal that might be attractive
precisely because it carries a known factor exposure rather than despite it.

\answerlines{6}

\section*{References}

\begin{enumerate}
  \item Fama and French (1993), \textit{Common Risk Factors in the Returns on Stocks and Bonds}.
  \item Grinold and Kahn, \textit{Active Portfolio Management}.
  \item Cochrane, \textit{Asset Pricing}.
  \item Lopez de Prado (2018), \textit{Advances in Financial Machine Learning}.
\end{enumerate}

\end{document}
